\documentclass[11pt]{amsart}
\usepackage{amssymb,booktabs}
\usepackage[hidelinks,pdftitle={A kissing configuration of 200540 points in 27 dimensions},pdfauthor={Alexey Kravatskiy}]{hyperref}

\newtheorem{theorem}{Theorem}
\newtheorem*{theorem*}{Theorem}
\newtheorem{lemma}[theorem]{Lemma}
\newtheorem{proposition}[theorem]{Proposition}
\theoremstyle{remark}
\newtheorem{remark}[theorem]{Remark}

% a symbol, not an operator: \DeclareMathOperator would make it a \mathop,
% adding operator spacing and detaching sub- and superscripts
\newcommand{\Leech}{\Lambda_{24}}
\newcommand{\ip}[2]{\langle #1,#2\rangle}
\newcommand{\wh}{\widehat{w}}

\title{A kissing configuration of $200540$ points in $\mathbb{R}^{27}$}
\author{Alexey Kravatskiy}
\address{MIRIAI (Moscow Independent Research Institute\newline
of Artificial Intelligence), Moscow, Russia}
\email{kravatskii.a@miriai.org}
\date{19 August 2026}
\subjclass[2020]{Primary 52C17; Secondary 11H31, 94B25}
\keywords{Kissing number, Leech lattice, Golay code, root system,
exact verification}

% pages end where their content ends; amsart's default \flushbottom stretches
% the vertical glue on the one page whose last display will not fit
\raggedbottom

\begin{document}

\begin{abstract}
We show that the kissing number in $27$ dimensions satisfies $K(27)\ge 200540$,
improving the record of $200044$ by $496$. The record configurations in
dimensions $25$ through $31$ all have a common form, and the number of points
in such a configuration is governed by a single integer $m$. We prove that
$m\le\lfloor 2K(k)/3\rfloor$ in dimension $24+k$, under hypotheses that all
seven published configurations satisfy, and find that all seven attain this
bound except in dimension $27$, where $m$ falls one short. Closing that gap
needs no new coordinates, only a regrouping of the configuration already
published,
and it adds $496$ points. For $k=5,6,7$ the value of $K(k)$ is not known, and
the comparison there is against the best lower bound known for it. Every
assertion made here about the configurations is verified in exact arithmetic
over $\mathbb Z[\sqrt2,\sqrt3,\sqrt6]$, the ring in which the published
coordinates lie.
\end{abstract}

\maketitle

The kissing number $K(n)$ is the largest number of unit vectors in
$\mathbb{R}^{n}$ whose pairwise inner products are all at most $\tfrac12$. We
prove the following, improving the record of $200044$ held by the configuration
of \cite{ma2025} by $496$.

\begin{theorem*}
$K(27)\ge 200540$.
\end{theorem*}

This is Theorem~\ref{thm:main} below. Section~\ref{sec:family} describes the
family of constructions the record configurations in dimensions $25$ through
$31$ all belong to, Section~\ref{sec:bound} bounds the quantity that governs
their size, Section~\ref{sec:dim27} gives the configuration, and
Section~\ref{sec:verification} verifies it.

\section{The construction family}\label{sec:family}

Write $\mathbb{R}^{24+k}=\mathbb{R}^{24}\oplus\mathbb{R}^{k}$ and split a unit
vector as $p=(v,w)$. Let $X\subset S^{23}$ denote the $196560$ minimal vectors
of the Leech lattice $\Leech$, normalised to unit length; recall that
$\ip{x}{x'}\in\{0,\pm\tfrac14,\pm\tfrac12,\pm1\}$ for $x,x'\in X$ \cite{cs};
we re-verified this value set in exact arithmetic for the copy of $X$ recovered
from \cite{cohn-data}.

Those configurations all have the following form, which we extracted directly
from the published coordinates \cite{cohn-data} for every $k=1,\dots,7$:

\begin{itemize}
\item an \emph{equator} $\{(x,0):x\in X\setminus D\}$, for a set $D\subset X$
      of \emph{deleted} vectors;
\item \emph{caps} $\bigl(\sqrt{\tfrac23}\,u,\ \sqrt{\tfrac13}\,\wh\bigr)$, in
      which $u\in X$ is the cap's \emph{head} and the unit vector
      $\wh\in\mathbb{R}^{k}$ its \emph{direction};
\item \emph{axis points} $(0,\omega)$, the $\omega$ forming a kissing
      configuration in $\mathbb{R}^{k}$: at most $K(k)$ of them, further
      constrained by \eqref{eq:A} below.
\end{itemize}

Let $W\subset S^{k-1}$ be the set of directions that occur and, for $\wh\in W$,
let $U_{\wh}\subset X$ be the set of heads occurring with $\wh$, the
\emph{class} of $\wh$. The constraints are as follows; all are immediate from
$\ip{p}{p'}\le\tfrac12$.

\begin{lemma}\label{lem:constraints}
The above configuration is a kissing configuration if and only if
\begin{align}
&\text{(E)} && u\in D \text{ for every } u\in\textstyle\bigcup_{\wh}U_{\wh},
   \label{eq:E}\\
&\text{(C)} && \tfrac23\ip{u}{u'}+\tfrac13\ip{\wh}{\wh{}'}\le\tfrac12
   \quad\text{for caps } (u,\wh)\ne(u',\wh{}'),\label{eq:C}\\
&\text{(A)} && \ip{\omega}{\wh}\le\tfrac{\sqrt3}{2}\quad\text{for all
   axis points }\omega\text{ and }\wh\in W.\label{eq:A}
\end{align}
\end{lemma}

\begin{proof}
An equatorial point $(x,0)$ and a cap $(\sqrt{2/3}\,u,\sqrt{1/3}\,\wh)$ have
inner product $\sqrt{2/3}\,\ip{x}{u}$, which is at most
$\sqrt{2/3}\cdot\tfrac12<\tfrac12$ when $x\ne u$, and equals
$\sqrt{2/3}>\tfrac12$ when $x=u$; this gives \eqref{eq:E}. Two caps give
\eqref{eq:C} directly, and an axis point and a cap give
$\sqrt{1/3}\,\ip{\omega}{\wh}\le\tfrac12$, which is \eqref{eq:A}. Equatorial
pairs and axis pairs are inherited from $X$ and from the kissing configuration
in $\mathbb{R}^{k}$, and $(x,0)\perp(0,\omega)$.
\end{proof}

Condition \eqref{eq:E} is the essential point: the head of a cap must be
deleted from the equator. Only then may the cap sit at
$|v|=\sqrt{2/3}>\tfrac12$.

\begin{remark}\label{rem:k1}
Condition \eqref{eq:A} admits no axis point at all when $k=1$: there
$W=\{\pm1\}$, so either choice of $\omega=\pm1$ has $\ip{\omega}{\wh}=1>
\tfrac{\sqrt3}{2}$ for one of the two cap directions. This is why the $k=1$ row
of Table~\ref{tab:m} carries no axis points where every other row carries
$K(k)$ of them.
\end{remark}

Specialising \eqref{eq:C} to $\wh{}'=\wh$ gives $\ip{u}{u'}\le\tfrac14$: every
class has pairwise cosines at most $\tfrac14$. Write
$\mathcal S=\{U_{\wh}:\wh\in W\}$ for the set of distinct classes.

\begin{proposition}\label{prop:count}
If the classes in $\mathcal S$ are pairwise disjoint and all have size $N$, and
$D=\bigcup_{\wh}U_{\wh}$, then the configuration has
\begin{multline*}
   \bigl(196560-|\mathcal S|\,N\bigr)\;+\;|W|\,N\;+\;(\text{\# axis points})\\
   {}=\;196560+\bigl(|W|-|\mathcal S|\bigr)N+(\text{\# axis points})
\end{multline*}
points. In particular the net gain is governed by $m:=|W|-|\mathcal S|$.
\end{proposition}

Each direction contributes $N$ points, while each \emph{distinct} class costs
$N$ deletions from the equator; a class reused on several directions is paid for
only once.

\section{The bound on \texorpdfstring{$m$}{m}}\label{sec:bound}

\begin{lemma}\label{lem:group}
Let $\wh\ne\wh{}'$ in $W$.
\begin{enumerate}
\item If $U_{\wh}\cap U_{\wh'}\ne\emptyset$ then $\ip{\wh}{\wh{}'}\le-\tfrac12$.
\item If there are $u\in U_{\wh}$ and $u'\in U_{\wh'}$ with
      $\ip{u}{u'}=\tfrac12$, then $\ip{\wh}{\wh{}'}\le\tfrac12$.
\end{enumerate}
\end{lemma}

\begin{proof}
(i) Taking $u=u'$ in the intersection in \eqref{eq:C} gives
$\tfrac23+\tfrac13\ip{\wh}{\wh{}'}\le\tfrac12$.
(ii) Applying \eqref{eq:C} to that pair gives
$\tfrac13+\tfrac13\ip{\wh}{\wh{}'}\le\tfrac12$.
\end{proof}

In (ii) it is essential that $u$ and $u'$ lie in \emph{different} classes: all
cosines within a class are at most $\tfrac14$, so a pair at cosine $\tfrac12$
inside a single class does not exist, and the hypothesis read that way would be
unsatisfiable.

\begin{theorem}\label{thm:bound}
Let a configuration be as in Proposition~\ref{prop:count}, so that the classes
$U_{\wh}$ are pairwise equal or disjoint. Then
\begin{enumerate}
\item $m=|W|-|\mathcal S|\le|W|-\lceil |W|/3\rceil=\lfloor 2|W|/3\rfloor$;
\item if in addition any two \emph{distinct} classes $U_{\wh}\ne U_{\wh'}$
      contain a cross pair $u\in U_{\wh}$, $u'\in U_{\wh'}$ with
      $\ip{u}{u'}=\tfrac12$, then $|W|\le K(k)$ and hence
      $m\le\lfloor 2K(k)/3\rfloor$.
\end{enumerate}
\end{theorem}

\begin{proof}
(i) Group the elements of $W$ according to which class they carry; this is a
partition because the classes are pairwise equal or disjoint. By
Lemma~\ref{lem:group}(i) the members of a group are pairwise at cosine at most
$-\tfrac12$, so if a group has $t$ members then
$0\le\bigl|\sum \wh\bigr|^{2}\le t+2\binom{t}{2}\bigl(-\tfrac12\bigr)
 =t-\tfrac{t(t-1)}{2}$,
whence $t\le3$. (This uses no property of $\mathbb{R}^{k}$ beyond the inner
product, so it holds in every dimension.) Therefore
$|\mathcal S|\ge\lceil |W|/3\rceil$.

(ii) Two directions carrying the same class are at cosine at most $-\tfrac12$
by Lemma~\ref{lem:group}(i); two carrying different classes are at cosine at
most $\tfrac12$ by Lemma~\ref{lem:group}(ii). So $W$ is a set of unit vectors
in $\mathbb{R}^{k}$ that are pairwise at angle at least $60^\circ$, i.e.\
$|W|\le K(k)$. Since $n\mapsto\lfloor 2n/3\rfloor$ is nondecreasing, (i) gives
the stated bound.
\end{proof}

\begin{remark}\label{rem:hyp}
Neither hypothesis can be dropped. Without the first, the classes carried by the
elements of $W$ need not be equal or disjoint, so the grouping used in the
proof of (i) is not a partition. Without the second, two disjoint classes all of
whose cross cosines are at most $0$ impose, through \eqref{eq:C}, only
$\ip{\wh}{\wh{}'}\le\tfrac32$, which is no constraint; $|W|$ could then exceed
$K(k)$. Both hypotheses do hold for all seven published configurations, as
recorded in Table~\ref{tab:m}.
\end{remark}

\begin{table}[t]
\centering
{\small\setlength{\tabcolsep}{4.5pt}
\begin{tabular}{ccccccccr}
\toprule
$k$ & $24+k$ & $W$ & $|W|$ & $|\mathcal S|$ & $m$
    & $\lfloor 2K(k)/3\rfloor$ & $K(k)$ & points\\
\midrule
\multicolumn{9}{l}{\itshape $K(k)$ known exactly}\\
$1$ & $25$ & $A_1$ & $2$   & $1$  & $1$  & $1$  & $2$  & $197056$\\
$2$ & $26$ & $A_2$ & $6$   & $2$  & $4$  & $4$  & $6$  & $198550$\\
$3$ & $27$ & $A_3$ & $12$  & $5$  & $\mathbf{7}$ & $\mathbf{8}$ & $12$ & $200044$\\
$4$ & $28$ & $D_4$ & $24$  & $8$  & $16$ & $16$ & $24$ & $204520$\\
\midrule
\multicolumn{9}{l}{\itshape $K(k)$ known only within the range shown}\\
$5$ & $29$ & $D_5$ & $40$  & $14$ & $26$ & $26$--$29$ & $40$--$44$   & $209496$\\
$6$ & $30$ & $E_6$ & $72$  & $24$ & $48$ & $48$--$51$ & $72$--$77$   & $220440$\\
$7$ & $31$ & $E_7$ & $126$ & $42$ & $84$ & $84$--$89$ & $126$--$134$ & $238350$\\
\bottomrule
\end{tabular}}
\caption{The bound of Theorem~\ref{thm:bound}(ii) against the values of $m$ in
the record configurations of \cite{ma2025,cohn-data}, all seven computed from
the published coordinates in exact arithmetic. The two middle columns are the
comparison, and they agree in every row but $k=3$, where $m=7$ falls one short
of $8$ --- the slack this note removes. In every row $W$ is closed under
negation and has the cosine distribution of the roots of the named root system,
so $|W|$ equals $K(k)$ in the upper block and the best known lower bound for
$K(k)$ in the lower one, where the comparison is accordingly conditional on the
configurations named in the $W$ column. Throughout, the classes are pairwise
disjoint with $N=496$; the axis points number $K(k)$ for $k\ge2$ and none for
$k=1$ (Remark~\ref{rem:k1}); and both hypotheses of Theorem~\ref{thm:bound}
hold, the maximum cross cosine over all $1267$ pairs of distinct classes
arising at $k=1,\dots,7$ being exactly $\tfrac12$.}
\label{tab:m}
\end{table}

\begin{remark}\label{rem:table}
The kissing number is known exactly only in dimensions $1$, $2$, $3$, $4$, $8$
and $24$; that is what separates the two blocks of Table~\ref{tab:m}. In the
lower block the entries $40$, $72$, $126$ are the best known lower bounds for
$K(5)$, $K(6)$, $K(7)$, each realised by the root system named in the same row,
and the best known upper bounds are $44$ and $134$ \cite{mv2010} and $77$
\cite{dllk2024}, as recorded in \cite{cohn-table}. Since $n\mapsto\lfloor2n/3\rfloor$ is
nondecreasing, the bound there is understated: were $K$ to equal its upper
bound, the entries would read $29$, $51$, $89$ and $m$ would fall short by $3$,
$3$, $5$. So in that block ``$m$ meets the bound'' asserts only that it meets
the value computed from the lower end of the range.

The distinction is not bookkeeping. What Theorem~\ref{thm:bound}(i) constrains
is $|\mathcal S|\ge\lceil|W|/3\rceil$, and six of the seven rows already sit at
that minimum: $|\mathcal S|=1,2,\mathbf{5},8,14,24,42$ against
$\lceil|W|/3\rceil=1,2,\mathbf{4},8,14,24,42$. Only $k=3$ spends a class it need
not, and that is the whole of the improvement below. For $k=5,6,7$ a larger $m$
therefore cannot come from regrouping a given $W$; it would need a larger $W$,
that is, a kissing configuration in $\mathbb{R}^{k}$ beating $D_5$, $E_6$ or
$E_7$ --- the open problem that put those rows in the lower block to begin with.
Rows $k=1,2,3,4$ use the proven values $K=2,6,12,24$, so the slack at $k=3$ is
unconditional.
\end{remark}

\section{Dimension \texorpdfstring{$27$}{27}}\label{sec:dim27}

We extracted the published $200044$-point configuration in $\mathbb{R}^{27}$
from \cite{cohn-data} and computed its structure. It consists of $194080$
equatorial points, $5952$ caps with $|w|=1/\sqrt3$ exactly, and $12$ axis
points. The caps lie on $|W|=12$ directions whose pairwise cosines lie in
$\{0,\pm\tfrac12,\pm1\}$, so $W$ is a cuboctahedron --- the $A_3$ root system of
Table~\ref{tab:m} --- and $|W|=K(3)=12$ is maximal. Each class has $N=496$
elements and there are $|\mathcal S|=5$ of them, pairwise disjoint, whose
$5\cdot496=2480$ vectors are exactly the $196560-194080=2480$ deleted from the
equator. Thus $m=12-5=7$, in agreement with Proposition~\ref{prop:count}:
$196560+7\cdot496+12=200044$.

The grouping of the twelve directions is $\{3,3,2,2,2\}$, with pairwise cosine
$-\tfrac12$ inside each triple and $-1$ inside each pair. By the proof of
Theorem~\ref{thm:bound}(i) a group has at most three members, so the minimum
possible number of groups is $\lceil 12/3\rceil=4$. We verified by exhaustive
search over the $\binom{12}{3}$ triples and $\binom{12}{4}$ quadruples that the
twelve directions contain exactly eight triples that are pairwise at cosine
$-\tfrac12$ (and no set of four with all cosines $\le-\tfrac12$) and that
these eight triples admit exactly two partitions of $W$ into four of them,
namely
\[
\begin{aligned}
   P_1:\quad&\{\wh_0,\wh_7,\wh_{10}\},&&\{\wh_1,\wh_6,\wh_9\},
   &&\{\wh_2,\wh_3,\wh_{11}\},&&\{\wh_4,\wh_5,\wh_8\};\\
   P_2:\quad&\{\wh_0,\wh_8,\wh_9\},   &&\{\wh_1,\wh_4,\wh_{11}\},
   &&\{\wh_2,\wh_5,\wh_{10}\},&&\{\wh_3,\wh_6,\wh_7\}.
\end{aligned}
\]
Here $\wh_0,\dots,\wh_{11}$ are the twelve directions in lexicographic order of
their coordinate triples; the statement is of course independent of the
labelling. Each of $P_1$ and $P_2$ reuses one of the two triples already present
in the published grouping. Assigning any four of the five existing classes to
the four triples of either partition yields $|\mathcal S|=4$, and hence $m=8$.

\begin{theorem}\label{thm:main}
$K(27)\ge 200540$.
\end{theorem}

\begin{proof}
Let $S_1,\dots,S_4$ be any four of the five classes. Take
$D=S_1\cup S_2\cup S_3\cup S_4$, of size $4\cdot496=1984$; the equator
$X\setminus D$ has $194576$ points. Place $S_i$ on the $i$-th triple of $P_1$
(or of $P_2$; the argument is identical), giving $12\cdot496=5952$ caps, and
retain the $12$ axis points. Equatorial
pairs are inherited from $X$ and axis pairs from the cuboctahedron of axis
points; $(x,0)\perp(0,\omega)$. Condition \eqref{eq:E} holds by construction,
since every cap head lies in $D$ and the equator is $X\setminus D$. Condition
\eqref{eq:A} holds because $W$ and the axis points are unchanged from the
published configuration:
$\max_{\omega,\wh}\ip{\omega}{\wh}=\tfrac12+\tfrac{\sqrt2}{4}=0.8535\ldots
<\tfrac{\sqrt3}{2}=0.8660\ldots$, equivalently the largest inner product between
a cap and an axis point is
$\tfrac{\sqrt3}{6}+\tfrac{\sqrt6}{12}=0.4927\ldots<\tfrac12$. For \eqref{eq:C}
there are three cases, using that distinct elements of $X$ have
$\ip{u}{u'}\le\tfrac12$ and that the four classes are pairwise disjoint, so that
$u=u'$ can occur only within one class:
\begin{itemize}
\item $\wh=\wh{}'$: then $u\ne u'$ lie in one class, so $\ip{u}{u'}\le\tfrac14$
      and $\tfrac23\cdot\tfrac14+\tfrac13=\tfrac12$;
\item $\wh\ne\wh{}'$ in the same triple: $\ip{\wh}{\wh{}'}=-\tfrac12$ and
      $\ip{u}{u'}\le1$, giving $\tfrac23-\tfrac16=\tfrac12$;
\item $\wh,\wh{}'$ in different triples: they carry different classes, which are
      disjoint, so $u\ne u'$ and $\ip{u}{u'}\le\tfrac12$; and
      $\ip{\wh}{\wh{}'}\le\tfrac12$ for any two distinct elements of a
      cuboctahedron. Hence
      $\tfrac23\cdot\tfrac12+\tfrac13\cdot\tfrac12=\tfrac12$, with equality
      exactly when both bounds are met.
\end{itemize}
The last case covers the remaining cuboctahedron cosines $0,-\tfrac12,-1$ across
triples as well, with strictly smaller values. Finally, an equatorial point and
a cap have inner product $\sqrt{2/3}\,\ip{x}{u}\le\sqrt{2/3}\cdot\tfrac12
=0.408\ldots$, since $x\ne u$ by \eqref{eq:E}. The total is
$194576+5952+12=200540$.
\end{proof}

By Theorem~\ref{thm:bound} this value of $m$ is optimal within the family:
$|W|\le K(3)=12$ and each group has at most three members, whence
$m\le\lfloor 2\cdot12/3\rfloor=8$. Since $K(3)=12$ is a theorem
\cite{svdw,cs}, this optimality statement is unconditional, in contrast to the
rows $k=5,6,7$ of Table~\ref{tab:m}. The configuration just built satisfies the
hypotheses of Theorem~\ref{thm:bound} as well as the published one does: its
four classes are four of the five of \cite{ma2025}, so they are still pairwise
disjoint of size $496$, and by Remark~\ref{rem:hyp} any two of them contain a
cross pair at cosine $\tfrac12$.

\begin{remark}\label{rem:N}
Theorem~\ref{thm:bound} bounds $m$, not the number of points, which by
Proposition~\ref{prop:count} is $196560+mN+(\text{\# axis points})$ and so is
linear in $N$. We prove nothing about $N$: we do not know the largest $N$ for
which $X$ contains $\lceil|W|/3\rceil$ pairwise disjoint $N$-element subsets of
pairwise cosine at most $\tfrac14$, and indeed we do not bound the size of even
one such subset. All seven published configurations use $N=496$ and the
improvement here keeps it, so Table~\ref{tab:m} compares configurations with the
same $N$; a larger admissible $N$ would give more points by the same counting.
Finding one, or ruling it out, is the natural next question for this family.
\end{remark}

\section{Verification}\label{sec:verification}

The configuration is available as \texttt{data/dimension27\_200540.txt} in the
accompanying verification package \cite{repo}, written in the format of
\cite{cohn-data}. The package proves the theorem twice, by different routes,
with no floating-point arithmetic and no tolerance at any step.

\texttt{scripts/verify\_configuration.py} works from the coordinate file alone,
needs no third-party packages, and takes some fifteen seconds on a laptop; it is
described below. \texttt{scripts/verify\_exhaustive.py} assumes nothing about the
geometry and compares all $\binom{200540}{2}=20\,108\,045\,530$ pairs directly.
Both also certify the published $200044$-point configuration of \cite{ma2025}
unchanged, which is a check on the verifiers as much as on the configuration.

Every coordinate in the dimension-$27$ block of \cite{cohn-data} lies in
$\mathbb{Z}[\sqrt2,\sqrt3,\sqrt6]$; the block uses $45$ distinct coordinate
tokens, all of that form. We represent the ring by integer $4$-tuples in the
basis $(1,\sqrt2,\sqrt3,\sqrt6)$, so ring arithmetic is integer arithmetic; the
comparisons that integers alone do not settle are settled by certified rational
enclosures of the radicals.

\begin{proposition}\label{prop:reduction}
After an exact rational rescaling, every point of the block in dimension $27$
takes one of three forms
\[
   \frac{(y,0)}{\sqrt{32}},\qquad
   \frac{(3y,t)}{\sqrt{432}},\qquad
   \frac{(0,a)}{\sqrt{48}},
\]
for the equatorial, cap and axis points respectively, where in the first two
$y\in\mathbb{Z}^{24}$ with $y\cdot y=32$, and where $t$ and $a$ lie in the ring
with $t\cdot t=144$ and $a\cdot a=48$. Consequently each pairwise inner product
is $\le\tfrac12$ precisely when the corresponding condition below holds:
\[
\begin{array}{lll@{\quad}l}
   \text{equator--equator} & \dfrac{y\cdot y'}{32}
      &\le\tfrac12 &\Longleftrightarrow\ y\cdot y'\le16,\\[8pt]
   \text{equator--cap}     & \dfrac{y\cdot y'}{16\sqrt6}
      &\le\tfrac12 &\Longleftrightarrow\ y\cdot y'\le\lfloor8\sqrt6\rfloor=19,\\[8pt]
   \text{cap--cap}         & \dfrac{9\,y\cdot y'+t\cdot t'}{432}
      &\le\tfrac12 &\Longleftrightarrow\ 9\,y\cdot y'\le216-t\cdot t',\\[8pt]
   \text{equator--axis}    & 0 && \text{always},\\[4pt]
   \text{cap--axis}        & \dfrac{t\cdot a}{144}
      &\le\tfrac12 &\Longleftrightarrow\ t\cdot a\le72,\\[8pt]
   \text{axis--axis}       & \dfrac{a\cdot a'}{48}
      &\le\tfrac12 &\Longleftrightarrow\ a\cdot a'\le24.
\end{array}
\]
The equator--equator, equator--cap and cap--cap conditions are comparisons of
integers; the cap--axis and axis--axis conditions are $12\times12$ tables of
comparisons in the ring.
\end{proposition}

\begin{proof}
The classification into the three forms is exact and scale invariant: the
equatorial points are the rows with $w=0$, the axis points those with $v=0$, and
the caps those with $|w|^{2}/|p|^{2}=\tfrac13$ as an identity between integers.
Only then is a row rescaled, to squared norm $32$, $432$ or $48$ according to its
family. The block contains rows of nine distinct squared norms, namely $2$, $3$,
$8$, $12$, $27$, $32$, $48$, $108$ and $432$; the factors required are
$4,2,1$ for the equatorial rows of norm $2,8,32$, then $12,6,4,3,2,1$ for the
cap rows of norm $3,12,27,48,108,432$, and $4,2,1$ for the axis rows of norm
$3,12,48$. Every factor is a positive integer, so the rescaling stays inside the
ring, and a rescaled cap head is divisible by $3$. The six displayed equivalences are
then immediate on clearing denominators, using $\sqrt{32\cdot432}=48\sqrt6$ in
the second, $\sqrt{432\cdot48}=144$ in the fifth, and
$19^{2}=361\le384<400=20^{2}$ to evaluate $\lfloor8\sqrt6\rfloor$.
\end{proof}

\begin{remark}\label{rem:scale}
The scaling above is the only one used anywhere in this note or in the
verification package. Fixing a single convention matters: an earlier
draft of the package inherited the cap--axis threshold from a normalisation in
which axis rows had squared norm $432$, and so tested $t\cdot a\le216$ rather
than $t\cdot a\le72$. Since $|t\cdot a|\le\sqrt{144\cdot48}=48\sqrt3<216$ always,
that test could not fail, and the tightest condition in the configuration --
$t\cdot a=24\sqrt3+12\sqrt6=70.96\ldots$ against a limit of $72$ -- went
unchecked. The package now derives each threshold from the squared norms of the
two families involved, and also asserts of each that Cauchy--Schwarz permits a
value above it, so that a vacuous test is itself a failure.
\end{remark}

\begin{remark}
Where a script evaluates the integer dot products with a floating point matrix
product, it is exact. The rational coordinates are bounded by $12$, so each of
the $24$ products has absolute value at most $144$ and every partial sum is an
integer of absolute value at most $3456$; integers of that size are represented
exactly in IEEE~754 binary64, so no rounding occurs at any step, under any
summation order and with or without fused multiply--add. The scripts do not
assume this: each checks a sample of its products against exact integer
arithmetic, and derives the bound from the data rather than assuming it.
\end{remark}

The outcome is that all $200540$ points are distinct and of norm exactly $1$,
and that every one of the $20\,108\,045\,530$ pairs satisfies
$\ip{p}{q}\le\tfrac12$.

The fast script avoids the quadratic sweep by certifying, from the coordinates
alone, that the $196560$ heads are the minimal vectors of a copy of $\Leech$.
Call the multiset of absolute values of the nonzero coordinates of a vector its
\emph{shape}. The census by shape is $1104$ of shape $(\pm4^{2})$, $97152$ of
shape $(\pm2^{8})$ and $98304$ of shape $(\mp3,\pm1^{23})$, which is the census
of $\Leech$; the supports of those of shape $(\pm2^{8})$ are $759$ octads,
spanning a $[24,12]$
binary code with the weight enumerator
$1+759x^{8}+2576x^{12}+759x^{16}+x^{24}$ of the extended Golay code and equal to
its weight-$8$ words; and every head satisfies the congruences defining $\Leech$
with respect to that code. Two distinct minimal vectors then satisfy
$y\cdot y'\le16$, which disposes of the equatorial pairs. As a further check on
the input, every one of the $196560$ was confirmed to have the inner product
distribution of the Leech minimal vectors: the numbers of $x'$ with
$\ip{x}{x'}$ equal to $-1,-\tfrac12,-\tfrac14,0,\tfrac14,\tfrac12,1$ are
$1,4600,47104,93150,47104,4600,1$, summing to $196560$.

Equality is attained, and attained exactly. The pairs realising $\tfrac12$ are
the equatorial pairs at Leech cosine $\tfrac12$, the $48$ ordered axis pairs at
cosine $\tfrac12$, and three families of cap pairs:
\[
\begin{array}{lll@{\qquad}r}
   \ip{\wh}{\wh{}'}=1,        &\ip{u}{u'}=\tfrac14:
      &\tfrac23\cdot\tfrac14+\tfrac13\cdot1=\tfrac12,
      &900096,\\[3pt]
   \ip{\wh}{\wh{}'}=\tfrac12, &\ip{u}{u'}=\tfrac12:
      &\tfrac23\cdot\tfrac12+\tfrac13\cdot\tfrac12=\tfrac12,
      &277408,\\[3pt]
   \ip{\wh}{\wh{}'}=-\tfrac12,&u=u':
      &\tfrac23\cdot1+\tfrac13\cdot\bigl(-\tfrac12\bigr)=\tfrac12,
      &11904,
\end{array}
\]
where the counts are of ordered pairs. The third family is the one created by
the regrouping: it consists of the $4\cdot6\cdot496$ pairs of caps that share a
head and sit on two distinct directions of a single triple. Its presence is the
reason a triple, and no larger set, is the most that can share a class.

\begin{remark}
The classes $S_1,\dots,S_4$ are taken verbatim from the configuration of
\cite{ma2025}: every Leech vector, direction and axis point used here already
occurs there, and no coordinate was computed anew. The improvement is purely the
observation that partitioning the twelve directions into four triples, rather
than two triples and three antipodal pairs, makes the fifth class unnecessary
and returns the $496$ vectors it was deleting to the equator.
\end{remark}

\section*{Use of AI assistance}

The construction of Section~\ref{sec:dim27} was found by Claude Opus~5 (Max
effort), at the fortieth prompt of a session in which the author supplied the
problem, Cohn's table \cite{cohn-table} and data set \cite{cohn-data}; most of
the thirty-nine earlier prompts were some form of ``try harder''. The model
maintained throughout that it would not succeed: ``I have not produced a new
SOTA \dots\ continuing would be me generating variations faster than either of
us can check them.'' This note and the verification of
Section~\ref{sec:verification} were then written by the author with Claude Code
(Opus~5, Extra effort) and attacked by further instances of the same model as
adversarial reviewers.

\begin{thebibliography}{9}
\bibitem{cohn-data} H.~Cohn, \emph{Table of kissing number bounds}, data set,
DSpace@MIT, \url{https://hdl.handle.net/1721.1/153312}.
\bibitem{cohn-table} H.~Cohn, \emph{Kissing numbers},
\url{https://cohn.mit.edu/kissing-numbers/}.
\bibitem{cs} J.~H.~Conway and N.~J.~A.~Sloane, \emph{Sphere packings, lattices
and groups}, 3rd ed., Springer, New York, 1999.
\bibitem{dllk2024} D.~de~Laat, N.~Leijenhorst, and W.~H.~H.~de~Muinck Keizer,
\emph{Optimality and uniqueness of the $D_4$ root system}, preprint, 2024,
arXiv:2404.18794.
\bibitem{repo} A.~Kravatskiy, \emph{Verification package for a kissing
configuration of $200540$ points in $\mathbb{R}^{27}$}, software and data, 2026,
\url{https://github.com/alexlegeartis/KissingNumbers}.
\bibitem{ma2025} Chengdong Ma, Th\'{e}o Tao Zhaowei, Pengyu Li, Minghao Liu,
Haojun Chen, Zihao Mao, Bo Li, Yuan Cheng, Yuan Qi, and Yaodong Yang,
\emph{Finding kissing numbers with game-theoretic reinforcement learning},
preprint, 2025, arXiv:2511.13391.
\bibitem{mv2010} H.~D.~Mittelmann and F.~Vallentin, \emph{High-accuracy
semidefinite programming bounds for kissing numbers}, Experiment.\ Math.\
\textbf{19} (2010), 175--179.
\bibitem{svdw} K.~Sch\"{u}tte and B.~L.~van~der~Waerden, \emph{Das Problem der
dreizehn Kugeln}, Math.\ Ann.\ \textbf{125} (1953), 325--334.
\end{thebibliography}

\end{document}
